\uuid{GZwl}
\exo7id{4335}
\titre{exo7 4335}
\auteur{quercia}
\organisation{exo7}
\datecreate{2010-03-12}
\isIndication{false}
\isCorrection{true}
\chapitre{Intégration}
\sousChapitre{Intégrale de Riemann dépendant d'un paramètre}
\cpge{mpsi-mp2i.limites-continuite.c, pcsi.limites-continuite.c, ptsi.limites-continuite.c, mpsi-mp2i.integration.c, pcsi.integration.b, ptsi.integration.b, bcpst1.limites-continuite.c, bcpst1.integration-segment.a, bcpst1.integration-segment.b, mp-mpi.integration-intervalle.f, bcpst2.integrales-generalisees}
\module{Analyse}
\niveau{L2}
\difficulte{}

\contenu{
\texte{
Soit $f : {[a,b]} \to {\R^+}$ continue atteignant son maximum en un unique
point $c \in {]a,b[}$.
}
\begin{enumerate}
    \item \question{Soit $\mu > 0$ tel que $[c-\mu,c+\mu] \subset [a,b]$.
    Chercher $\lim_{n\to\infty} \biggl( \int_{t=a}^b f^n(t)\,d t\biggm/ \int_{t=c-\mu}^{c+\mu} f^n(t)\,d t\biggr)$.}
    \item \question{Soit $g : {[a,b]} \to {\R^+}$ continue.
    Chercher $\lim_{n\to\infty} \biggl( \int_{t=a}^b f^n(t)g(t)\,d t\biggm/ \int_{t=a}^b f^n(t)\,d t\biggr)$.}
\reponse{
$g(c)$.
}
\end{enumerate}
}
